\appendix
\section*{附录}
\setcounter{subsection}{0}
\renewcommand{\thesubsection}{附录 \Alph{subsection}}

\subsection{支撑材料清单}

\begin{center}\small
\renewcommand{\arraystretch}{1.2}
\begin{tabular}{C{4.2cm}C{10.4cm}}
\toprule
文件 & 说明 \\
\midrule
\texttt{result1.xlsx}--\texttt{result4.xlsx} & 四问的完整结果 \\
\texttt{drying\_model.py} & 四问求解程序：读入附件，生成 result1--4 及表 1--6 的数据（附录 C） \\
\texttt{analysis.py} & 分析程序：湿空气状态、两阶段与干燥曲线、能量一致模型、工艺情景、问题四收缩关系与数值检验，结果写为 CSV（附录 D） \\
\texttt{figures.py} & 由 \texttt{analysis.py} 的 CSV 绘制论文全部插图 \\
\texttt{exports/redo/*.csv} & \texttt{analysis.py} 的输出数据，正文第 2、6--12 节的数值均出自这些文件 \\
\texttt{附件1.xlsx}、\texttt{附件2.xlsx} & 题给的烘房温湿度与药材半径数据 \\
\bottomrule
\end{tabular}
\end{center}

\subsection{运行环境与复现方式}

程序在 Python 3.13 下运行，依赖 numpy 2.3、scipy 1.16、openpyxl 3.1、matplotlib 3.10。把附件放在 \texttt{附件/} 目录后依次执行：
\begin{enumerate}[label=(\arabic*),leftmargin=2.2em,itemsep=1pt]
\item \texttt{python drying\_model.py -{}-data 附件 -{}-out output}：生成 result1--4.xlsx 与表 1--6 的 CSV；
\item \texttt{python analysis.py -{}-data 附件 -{}-out exports/redo}：生成第 2、6--12 节所用的全部数据；
\item \texttt{python figures.py}：读取上述 CSV，生成全部插图。
\end{enumerate}

\subsection{四问求解程序 drying\_model.py}

\lstinputlisting{code/drying_model.py}

\subsection{分析程序 analysis.py}

\lstinputlisting{code/analysis.py}
